\section{Ethical Considerations}
\label{sec:ethics}
CoMem changes the cost and storage organization of an underlying language model;
it does not add a safety mechanism. It therefore inherits risks such as
hallucination, biased generation, unsafe completion, and unintended disclosure
of sensitive text selected from a memory store. Lower-cost access to long
histories could also scale benign applications and deliberate misuse alike.
Cached residual tensors may also expose source information through inversion or
membership-inference attacks; they should receive protections at least as strong
as the source text. We do not empirically evaluate these attacks, so the risk is
a deployment concern rather than a demonstrated vulnerability. Their large persistent volume and possible cross-tenant reuse
make authorization, isolation, and deletion operational requirements rather than
properties conferred by non-readable tensor storage. Deployments should encrypt and audit the persistent store,
enforce access control, filter or redact sensitive sources before Write, support
verifiable deletion, restrict retrieval to authorized material, and retain
application-appropriate output monitoring and human review. Although bounded reads reduce accelerator memory
and per-query computation, training and large-scale evaluation still consume
energy; we report the measured hardware and final training budget in
Appendix~\ref{app:reproducibility}.

We collect no new human-subject data and recruit no annotators. Our experiments
use previously released model, corpus, and benchmark artifacts; LoCoMo's released
conversations are generated by paired language-model agents rather than newly
collected private conversations. The submitted software archive contains code
only and excludes model weights, benchmark text, predictions, credentials, and
API responses.
